\documentclass[10pt,aspectratio=169]{beamer}
\input{../../_en_preambulo_beamer}% Comandos y macros estadísticas compartidas para las presentaciones Beamer en inglés (Metropolis)

% Conjuntos numéricos básicos
\newcommand{\R}{\mathbb{R}}
\newcommand{\N}{\mathbb{N}}
\newcommand{\Q}{\mathbb{Q}}
\newcommand{\Z}{\mathbb{Z}}
\newcommand{\set}[1]{\left\{ #1 \right\}}
\newcommand{\abs}[1]{\left|#1\right|}
\newcommand{\norm}[1]{\left\| #1 \right\|}

% Comandos estadísticos y probabilísticos del curso
\newcommand{\Var}{\operatorname{Var}}
\newcommand{\cov}{\operatorname{Cov}}
\newcommand{\corr}{\rho}
\newcommand{\comb}[2]{\begin{pmatrix} #1 \\ #2 \end{pmatrix}}
\newcommand{\s}{\sigma}
\newcommand{\particion}{\mathfrak{P}}
\newcommand{\card}[1]{\#\left( #1 \right)}
\newcommand{\seq}[3]{\set{#1}_{#2}^{#3}}
\newcommand{\E}{\mathbb{E}}
\newcommand{\Prob}{\mathbb{P}}

% Entornos pedagógicos y cajas en inglés académico natural (soportando tanto nombres en inglés como en español)
\newenvironment{discussion}{%
  \begin{alertblock}{Discussion Question}%
}{%
  \end{alertblock}%
}
\newenvironment{discusion}{%
  \begin{alertblock}{Discussion Question}%
}{%
  \end{alertblock}%
}

\newenvironment{reflection}{%
  \begin{alertblock}{Classroom Reflection}%
}{%
  \end{alertblock}%
}
\newenvironment{reflexion}{%
  \begin{alertblock}{Classroom Reflection}%
}{%
  \end{alertblock}%
}

\newenvironment{paradox}{%
  \begin{alertblock}{Exploratory Paradox}%
}{%
  \end{alertblock}%
}
\newenvironment{paradoja}{%
  \begin{alertblock}{Exploratory Paradox}%
}{%
  \end{alertblock}%
}

\newenvironment{motivation}{%
  \begin{exampleblock}{Motivation: Data Science in Practice}%
}{%
  \end{exampleblock}%
}
\newenvironment{motivacion}{%
  \begin{exampleblock}{Motivation: Data Science in Practice}%
}{%
  \end{exampleblock}%
}

\newenvironment{quickchallenge}{%
  \begin{block}{Quick Team Challenge}%
}{%
  \end{block}%
}
\newenvironment{retoclase}{%
  \begin{block}{Quick Team Challenge}%
}{%
  \end{block}%
}

\title[Blocks and Latin squares]{Section 08.06: Blocks, Latin, and Graeco-Latin Squares}\subtitle{Controlling variability sources}\author[J. Castillo Colmenares]{Juliho Castillo Colmenares}\institute{Tecnológico de Monterrey}\date{}
\begin{document}
\begin{frame}[plain]\titlepage\end{frame}
\begin{frame}{Roadmap}\small\begin{enumerate}\item RCBD and additive model.\item Latin square.\item Graeco-Latin square.\end{enumerate}\end{frame}
\begin{frame}{Source and scope}\small This section introduces designs that control one, two, or three known heterogeneity sources in addition to treatment.\begin{block}{Guiding question}How can noise be reduced without unnecessary runs?\end{block}\end{frame}
\begin{frame}{Why block}\small When units differ in a known way, a completely randomized design mixes that variation with error. Blocks separate it.\end{frame}
\begin{frame}{RCBD}\small A Randomized Complete Block Design has $b$ homogeneous blocks and $k$ treatments; every treatment appears once per block, $N=kb$.\end{frame}
\begin{frame}{RCBD additive model}\small\[Y_{ij}=\mu+\tau_i+\beta_j+\epsilon_{ij}.\]\begin{itemize}\item $\tau_i$: treatment.\item $\beta_j$: block.\item $\epsilon_{ij}$: error.\end{itemize}\end{frame}
\begin{frame}{RCBD partition}\small\[\mathrm{SCT}=\mathrm{SCTR}+\mathrm{SCB}+\mathrm{SCE}.\]Block variation is removed before assessing treatments.\end{frame}
\begin{frame}{Latin square}\small A Latin square of order $n$ is an $n\times n$ array where each treatment appears exactly once in every row and column.\end{frame}
\begin{frame}{Two blocking sources}\small Rows and columns represent two crossed control factors. The design uses $n^2$ runs rather than the $n^3$ required by a complete three-factor design.\end{frame}
\begin{frame}{Latin model}\small\[Y_{ijk}=\mu+\rho_i+\gamma_j+\tau_k+\epsilon_{ijk}.\]\end{frame}
\begin{frame}{Latin partition}\small\[\mathrm{SCT}=\mathrm{SC}_{\mathrm{rows}}+\mathrm{SC}_{\mathrm{columns}}+\mathrm{SCTR}+\mathrm{SCE}.\]\end{frame}
\begin{frame}{Graeco-Latin square}\small Superimpose two orthogonal Latin squares: the Latin array represents the main treatment and Greek letters a fourth controlled source.\end{frame}
\begin{frame}{Existence condition}\small No orthogonal Graeco-Latin squares exist for order $6$; constructions exist for other orders $n\geq3$.\end{frame}
\begin{frame}{Graeco-Latin model}\small\[Y_{ijkl}=\mu+\rho_i+\gamma_j+\tau_k+\delta_l+\epsilon_{ijkl}.\]\end{frame}
\begin{frame}{Graeco-Latin partition}\small\[\mathrm{SCT}=\mathrm{SC}_{\mathrm{rows}}+\mathrm{SC}_{\mathrm{columns}}+\mathrm{SCTR}+\mathrm{SC}_{\mathrm{Greek}}+\mathrm{SCE}.\]\end{frame}
\begin{frame}{Application: RCBD}\small To compare treatments on servers from different generations, use generation as a block and randomize treatments within it.\end{frame}
\begin{frame}{Application: Latin square}\small If shift and machine are crossed noise sources, use rows and columns to control them while each treatment appears once in each direction.\end{frame}
\begin{frame}{Application: means}\small Compare treatment means after removing systematic contributions from blocks, rows, or columns.\end{frame}
\begin{frame}{Interpretation}\small An effective block reduces CME and can increase the power of the treatment $F$ statistic.\end{frame}
\begin{frame}{Communication}\small Report structure, constraints, sums of squares, degrees of freedom, and the treatment hypothesis; blocks are controls, not necessarily targets.\end{frame}
\begin{frame}{Synthesis}\small\begin{itemize}\item RCBD controls one source.\item Latin squares control two.\item Graeco-Latin squares control three.\item The goal is smaller residual error.\end{itemize}\end{frame}
\begin{frame}{Chapter connection}\small These designs extend ANOVA and prepare factorial designs with more than one treatment factor.\end{frame}
\end{document}
